
%% ========================================================================
\section{Methods}
\label{sec:methods}
%% ========================================================================

\subsection{Development Process and Timeline}

The repository and this paper were constructed using an AI-assisted development methodology. Kevin Kawchak served as the prompt creator, system architect, and provided edits to the generated paper, directing AI code generation through structured prompts documented in \texttt{prompts.md}.

The primary AI system used for repository generation and this paper was \textbf{Anthropic Claude Code Opus 4.6}. Peer review was conducted using \textbf{OpenAI ChatGPT 5.4 Thinking}, with review documents archived under the \texttt{peer-review/} directory.

\subsubsection{Repository Evolution}

The system was developed across four repositories in sequence. Table~\ref{tab:repo-timeline} summarizes the development timeline:

\begin{table}[H]
\centering
\caption{Repository development timeline and evolution.}
\label{tab:repo-timeline}
\scriptsize
\begin{tabularx}{\textwidth}{@{}L{4.7cm}L{1.6cm}R{0.7cm}L{5.2cm}@{}}
\toprule
\textbf{Repository} & \textbf{Date Range} & \textbf{Days} & \textbf{Focus} \\
\midrule
\texttt{physical-ai-oncology-trials}~\cite{kawchak2026physicalai} & Feb 2026 & 14 & USL scoring, patient instructions, single-site Physical AI framework \\
\texttt{pai-oncology-trial-fl}~\cite{kawchak2026paifl} & Feb--Mar 2026 & 10 & Federated learning, privacy modules, regulatory compliance \\
\texttt{mcp-pai-oncology-trials}~\cite{kawchak2026trialmcp} & Mar 2026 & 7 & MCP server proof of concept, 5-server architecture, 23 tools \\
\texttt{national-mcp-pai-oncology-trials}~\cite{kawchak2026national} & Mar 2026 & 5 & National standard, SDKs, conformance suite, governance \\
\bottomrule
\end{tabularx}
\end{table}

\subsubsection{Version History and AI Contributions}

The national repository progressed through the following versions, each driven by AI-assisted code generation:

\begin{itemize}[leftmargin=*]
\item \textbf{v0.5.0 (March 5, 2026):} Initial repository scaffold with 5 MCP server implementations, 13 JSON schemas, 8 conformance profiles, 44 unit tests, and deployment configurations. Generated by Claude Code Opus 4.6 from a structured prompt.

\item \textbf{v0.5.1 (March 6, 2026):} Added 34 integration adapters, 8 safety modules, interoperability testbed with 8 scenarios, benchmark suite, and expanded test coverage to 331 conformance tests plus 293 additional unit tests. This version addressed gaps identified by OpenAI ChatGPT 5.4 Thinking peer review.

\item \textbf{v0.5.2 (March 6, 2026):} Incorporated peer review feedback: full server implementations with transport, routing, and middleware; persistent storage adapters (SQLite, PostgreSQL); Docker deployment stack; and Kubernetes manifests. AI fixes were applied across 2 days.

\item \textbf{v1.0.0 (March 7, 2026):} Phase 5 release with Python and TypeScript SDKs, CLI tooling, code generation pipeline, stakeholder guides, operational documentation, and governance framework. Total: 381 files, 88 directories, approximately 69,800 lines.

\item \textbf{v1.0.1 (March 8, 2026):} Documentation accuracy corrections: badge updates, test count reconciliation, CI pipeline documentation alignment (14 jobs).

\item \textbf{v1.1.0 (March 9, 2026):} Initial paper, updated documentation, and release notes. Revised paper with formatting corrections: running header sizing, introduction restructuring, table column overlap fixes, and diagram pagination.
\end{itemize}

\subsubsection{Development Metrics}

Table~\ref{tab:dev-metrics} presents quantitative development metrics derived from release dates and version identifiers. Earlier v0.1.0 through v0.4.0 developments are also available in the repository. 

\begin{table}[H]
\centering
\caption{Development metrics for the national repository (v0.5.0 through v1.1.0).}
\label{tab:dev-metrics}
\small
\begin{tabular}{@{}lr@{}}
\toprule
\textbf{Metric} & \textbf{Value} \\
\midrule
Total span (v0.5.0 to v1.1.0) & 5 days \\
Days with an AI update & 5 days \\
Number of AI code generation sessions & 7 versions \\
Days of AI code recommendations & 4 days \\
Days of AI fixes (v0.5.2, v1.0.1) & 2 days \\
Total span across all 4 repositories & 36 days \\
\bottomrule
\end{tabular}
\end{table}

\subsection{Prompt Engineering Methodology}

Each version was generated from a structured prompt that specified the target architecture, file structure, and quality requirements. The prompts are archived in \texttt{prompts.md} and follow a consistent pattern:

\begin{enumerate}[leftmargin=*]
\item Define the target state (e.g., ``implement 5 MCP servers with 23 tools'').
\item Specify file paths and module structure.
\item List quality constraints (linting, testing, documentation).
\item Request CI/CD compatibility (ruff, pytest, GitHub Actions).
\end{enumerate}

This paper was generated from a prompt specifying the 20-page structure, LaTeX template (arxiv-style~\cite{arxivstyle}), section requirements, and formatting constraints. The prompt is archived in \texttt{prompts.md} under the v1.1.0 entry.

\subsection{Peer Review Process}

Peer review was conducted using OpenAI ChatGPT 5.4 Thinking, which analyzed the repository structure and provided recommendations documented in the \texttt{peer-review/} directory. Two recommendations by OpenAI (\texttt{openai-1-2026-03-06.md}, \texttt{openai-2-2026-03-06.md}), and five resulting new prompts by Claude (\texttt{claude-1-2026-03-06.md} through \texttt{claude-5-2026-03-06.md}) addressed gaps in:

\begin{itemize}[leftmargin=*]
\item Server implementation completeness (transport, routing, persistence)
\item Deployment readiness (Docker, Kubernetes, Helm)
\item Test coverage expansion (adversarial, security, interoperability)
\item SDK and developer tooling
\end{itemize}

A response document (\texttt{response-A-2026-03-06.md}) tracked how each recommendation was addressed in subsequent versions.

%% ========================================================================
\section{Results}
\label{sec:results}
%% ========================================================================

\subsection{Five-Server MCP Architecture}

The system implements five domain-specific MCP servers, each encapsulating a distinct responsibility within the oncology trial workflow. Together, they provide the complete protocol surface required for Physical AI integration with clinical trial systems.

\begin{table}[H]
\centering
\caption{Five MCP servers, their tools, and conformance levels.}
\label{tab:five-servers}
\footnotesize
\begin{tabularx}{\textwidth}{@{}L{3.2cm}C{0.8cm}C{1.0cm}L{6.8cm}@{}}
\toprule
\textbf{Server} & \textbf{Tools} & \textbf{Level} & \textbf{Primary Function} \\
\midrule
\texttt{trialmcp-authz} & 5 & 1 & Deny-by-default RBAC, SHA-256 token lifecycle, policy evaluation \\
\texttt{trialmcp-fhir} & 4 & 2 & FHIR R4 clinical data with HIPAA Safe Harbor de-identification \\
\texttt{trialmcp-dicom} & 4 & 3 & DICOM imaging queries, RECIST~\cite{eisenhauer2009recist} measurements \\
\texttt{trialmcp-ledger} & 5 & 1 & Hash-chained 21 CFR Part 11 audit trail \\
\texttt{trialmcp-provenance} & 5 & 4 & DAG-based data lineage, SHA-256 fingerprinting \\
\midrule
\textbf{Total} & \textbf{23} & & \\
\bottomrule
\end{tabularx}
\end{table}

\subsubsection{Authorization Server}

The authorization server (\texttt{servers/trialmcp\_authz/}) implements deny-by-default role-based access control (RBAC) as defined in ADR-007 (\texttt{docs/adr/007-deny-by-default-rbac.md}). It exposes five tools:

\begin{enumerate}[leftmargin=*]
\item \texttt{authz\_evaluate}: Evaluates whether a given role is permitted to invoke a specific tool.
\item \texttt{authz\_issue\_token}: Issues a scoped, time-limited token with SHA-256 hashing.
\item \texttt{authz\_validate\_token}: Validates a token's authenticity and expiration.
\item \texttt{authz\_revoke\_token}: Revokes a previously issued token.
\item \texttt{authz\_list\_policies}: Lists active policies for a given role or server.
\end{enumerate}

The server is implemented in \texttt{servers/trialmcp\_authz/server.py} (134 lines) with a policy engine (\texttt{policy\_engine.py}) and token store (\texttt{token\_store.py}).

\subsubsection{FHIR Clinical Data Server}

The FHIR server (\texttt{servers/trialmcp\_fhir/}) provides clinical data access with mandatory HIPAA Safe Harbor de-identification. The de-identification pipeline (\texttt{deid\_pipeline.py}) implements HMAC-SHA256 pseudonymization to maintain referential integrity across de-identified records. The four tools are:

\begin{itemize}[leftmargin=*]
\item \texttt{fhir\_read}: Read a single FHIR R4 resource by type and ID.
\item \texttt{fhir\_search}: Search FHIR resources with parameter filtering (limit: 100).
\item \texttt{fhir\_patient\_lookup}: Look up a patient by pseudonymized identifier.
\item \texttt{fhir\_study\_status}: Retrieve the status of a clinical study.
\end{itemize}

\subsubsection{Audit Ledger Server}

The ledger server (\texttt{servers/trialmcp\_ledger/}) implements a hash-chained immutable audit trail satisfying 21 CFR Part 11 requirements~\cite{fda2003part11}. The chain implementation in \texttt{chain.py} uses SHA-256 hashing with canonical JSON serialization. The following code excerpt shows the core hashing mechanism:

\begin{lstlisting}[caption={Hash-chained audit record computation from \texttt{servers/trialmcp\_ledger/chain.py}.},label={lst:chain}]
GENESIS_HASH = "0" * 64

def canonical_json(record: dict[str, Any]) -> str:
    filtered = {k: v for k, v in sorted(record.items())
                if k != "hash"}
    return json.dumps(filtered, sort_keys=True,
                      ensure_ascii=True)

def compute_audit_hash(record, prev_hash: str) -> str:
    canonical = canonical_json(record)
    payload = prev_hash + canonical
    return hashlib.sha256(
        payload.encode("utf-8")).hexdigest()
\end{lstlisting}

Each audit record contains: \texttt{audit\_id}, \texttt{timestamp}, \texttt{server}, \texttt{tool}, \texttt{caller}, \texttt{parameters}, \texttt{result\_summary}, \texttt{previous\_hash}, and \texttt{hash}. The chain is verified by recomputing each hash and comparing it against the stored value.

\subsubsection{Provenance Server}

The provenance server (\texttt{servers/trialmcp\_provenance/}) maintains a directed acyclic graph (DAG) of data lineage records, enabling both forward and backward lineage queries. The DAG implementation in \texttt{dag.py} supports five source types (\texttt{fhir}, \texttt{dicom}, \texttt{computed}, \texttt{external}, \texttt{federated}) and five action types (\texttt{read}, \texttt{transform}, \texttt{aggregate}, \texttt{transfer}, \texttt{derive}).

\subsubsection{DICOM Imaging Server}

The DICOM server (\texttt{servers/trialmcp\_dicom/}) provides imaging data access with role-based permissions and support for RECIST 1.1~\cite{eisenhauer2009recist} measurements critical to oncology trial response evaluation. The server exposes four tools: \texttt{dicom\_query} for searching imaging studies by patient or accession number, \texttt{dicom\_retrieve\_pointer} for obtaining secure references to image data without transferring pixel data, \texttt{dicom\_study\_metadata} for retrieving study-level DICOM metadata, and \texttt{dicom\_recist\_measurements} for querying tumor measurement data used in treatment response assessment.

The DICOM adapter (\texttt{dicom\_adapter.py}) supports integration with industry-standard PACS systems through the 9 DICOM integration adapters in \texttt{integrations/dicom/}, including dcm4chee, Orthanc, and DICOMweb connectors.

\subsubsection{Server Architecture Diagram}

The following text diagram shows the complete five-server deployment topology with data flow paths:

\begin{figure}[H]
\begin{verbatim}
  NATIONAL MCP FIVE-SERVER DEPLOYMENT
  +---------------------------------------------------+
  |                  SITE DEPLOYMENT                  |
  |                                                   |
  |    +-----------+  +-----------+  +-----------+    |
  |    | trialmcp  |  | trialmcp  |  | trialmcp  |    |
  |    | authz     |  | fhir      |  | dicom     |    |
  |    | (5 tools) |  | (4 tools) |  | (4 tools) |    |
  |    | RBAC,     |  | FHIR R4,  |  | DICOMweb, |    |
  |    | Tokens    |  | HIPAA     |  | RECIST    |    |
  |    +-----+-----+  +-----+-----+  +-----+-----+    |
  |          |              |              |          |
  |          +------+-------+------+-------+          |
  |                 |              |                  |
  |          +------v------+ +-----v-------+          |
  |          | trialmcp    | | trialmcp    |          |
  |          | ledger      | | provenance  |          |
  |          | (5 tools)   | | (5 tools)   |          |
  |          | SHA-256     | | DAG, W3C    |          |
  |          | 21 CFR 11   | | PROV        |          |
  |          +-------------+ +-------------+          |
  |                                                   |
  |  23 tools total across 5 servers                  |
  +---------------------------------------------------+
\end{verbatim}
\end{figure}

\subsubsection{Complete Tool Inventory}

Table~\ref{tab:tool-inventory} provides the complete inventory of all 23 tools with their input parameters and response types:

\begin{table}[H]
\centering
\caption{Complete 23-tool inventory across all five MCP servers.}
\label{tab:tool-inventory}
\footnotesize
\begin{tabularx}{\textwidth}{@{}L{4.2cm}L{3.5cm}L{4.1cm}@{}}
\toprule
\textbf{Tool} & \textbf{Key Parameters} & \textbf{Response Type} \\
\midrule
\texttt{authz\_evaluate} & role, tool & ALLOW/DENY decision \\
\texttt{authz\_issue\_token} & role, expires\_in & Token with SHA-256 hash \\
\texttt{authz\_validate\_token} & token\_hash & Valid/invalid status \\
\texttt{authz\_revoke\_token} & token\_hash & Revocation confirmation \\
\texttt{authz\_list\_policies} & role, server & Policy list \\
\midrule
\texttt{fhir\_read} & resource\_type, id & De-identified FHIR resource \\
\texttt{fhir\_search} & resource\_type, params & De-identified bundle \\
\texttt{fhir\_patient\_lookup} & pseudonym & De-identified patient \\
\texttt{fhir\_study\_status} & study\_id & Study status object \\
\midrule
\texttt{dicom\_query} & patient\_id, modality & Study result list \\
\texttt{dicom\_retrieve\_pointer} & study\_uid & Secure image pointer \\
\texttt{dicom\_study\_metadata} & study\_uid & DICOM metadata \\
\texttt{dicom\_recist\_measurements} & study\_uid & RECIST measurements \\
\midrule
\texttt{ledger\_append} & server, tool, caller & Audit record with hash \\
\texttt{ledger\_verify} & (none) & Chain validity status \\
\texttt{ledger\_query} & server, tool, caller & Filtered audit records \\
\texttt{ledger\_replay} & start\_index, count & Record sequence \\
\texttt{ledger\_export} & format & Full chain export \\
\midrule
\texttt{provenance\_record} & source\_id, action, actor & Provenance record \\
\texttt{provenance\_query\_forward} & record\_id, depth & Downstream lineage \\
\texttt{provenance\_query\_backward} & record\_id, depth & Upstream lineage \\
\texttt{provenance\_fingerprint} & data & SHA-256 fingerprint \\
\texttt{provenance\_lineage} & source\_id & Full lineage graph \\
\bottomrule
\end{tabularx}
\end{table}

\subsection{Conformance Levels and Profiles}

The standard defines five hierarchical conformance levels, each building upon the previous:

\begin{table}[H]
\centering
\caption{Five hierarchical conformance levels with required servers and use cases.}
\label{tab:conformance-levels}
\small
\begin{tabularx}{\textwidth}{@{}C{0.7cm}L{2.2cm}L{3.8cm}L{5cm}@{}}
\toprule
\textbf{Level} & \textbf{Name} & \textbf{Required Servers} & \textbf{Use Case} \\
\midrule
1 & Core & authz + ledger & Basic authorization and audit \\
2 & Clinical Read & Level 1 + fhir & FHIR clinical data access \\
3 & Imaging & Level 2 + dicom & DICOM imaging queries \\
4 & Federated Site & Level 3 + provenance & Multi-site data lineage \\
5 & Robot Procedure & All 5 servers & Full autonomous clinical workflow \\
\bottomrule
\end{tabularx}
\end{table}
